% chapters/12/section-04.tex
\section{Scene, Project, dan Penyimpanan}
\label{sec:scene-project}

\subsection{Scene sebagai Satu Level}

Sebuah \istilah{scene} adalah satu tingkat permainan, yaitu kumpulan benda
beserta seluruh pengaturannya. Sebuah permainan dapat tersusun atas beberapa
\istilah{scene}, misalnya satu \istilah{scene} untuk menu utama dan satu
\istilah{scene} untuk setiap tingkat permainan.

Pada Pygame, padanan \istilah{scene} adalah keadaan yang Anda susun sendiri di
dalam kode: daftar objek, posisi awalnya, dan pengaturannya. Pada Unity, seluruh
keadaan itu disimpan sebagai berkas, sehingga dapat dibuka, disunting, dan
dijalankan tanpa menulis kode.

\subsection{Struktur Berkas Project}

Setiap \istilah{project} Unity memiliki struktur folder yang tetap. Tiga di
antaranya perlu Anda ketahui.

\begin{description}
  \item[\code{Assets}] Memuat seluruh berkas yang Anda buat dan impor: scene,
        \istilah{script}, gambar, bunyi, dan aset lainnya. Folder inilah yang
        perlu Anda simpan pada repositori Git.
  \item[\code{Library}] Memuat berkas sementara hasil pengolahan Unity. Folder
        ini dapat dibangun ulang dan \textbf{tidak boleh} disimpan pada
        repositori, sebab ukurannya besar dan isinya berbeda pada setiap
        komputer.
  \item[\code{ProjectSettings}] Memuat pengaturan project, misalnya nama
        project, ukuran layar, dan \istilah{scene} mana yang dijalankan lebih
        dahulu.
\end{description}

\begin{peringatanbox}
Pada blok Python, berkas yang tidak perlu disimpan didaftarkan pada
\code{.gitignore}. Kebiasaan yang sama berlaku di sini, dan bahkan lebih
penting: folder \code{Library} pada project Unity dapat berukuran beberapa
gigabita. Daftarkan folder \code{Library}, \code{Temp}, \code{Obj}, dan
\code{Build} pada berkas \code{.gitignore}, sedangkan folder \code{Assets} dan
\code{ProjectSettings} harus disimpan.
\end{peringatanbox}

\subsection{Menyimpan Pekerjaan}

Unity menyimpan scene tidak secara otomatis, sehingga penyimpanan perlu
dilakukan dengan sengaja. Dua hal perlu dibedakan.

\begin{itemize}
  \item \textbf{Menyimpan scene} menyimpan susunan benda beserta pengaturannya.
        Perintahnya tersedia pada menu \istilah{File}.
  \item \textbf{Menyimpan \istilah{script}} menyimpan kode Anda. Berkas skrip
        disimpan oleh penyunting kode yang Anda pakai, bukan oleh Unity, dan
        Unity akan memuat ulang skripnya ketika Anda kembali ke jendelanya.
\end{itemize}

\begin{tipsbox}
Urutan yang menghemat waktu: simpan skrip lebih dahulu melalui penyunting kode,
baru simpan scene melalui Unity, lalu tekan tombol jalankan. Urutan itu
menghindarkan Anda dari menjalankan scene dengan skrip versi lama, suatu
keadaan yang membingungkan karena perubahan Anda seolah-olah tidak
berpengaruh.
\end{tipsbox}
